\section{Signs on every radius and parameter}\label{sec:signs}

Fix a real \(k\ge2.36\), and take the exact lists \(\ell_1,\ell_2\) from
Proposition~\ref{prop:exact-interpolation}.  Write \(p_i=p_{\ell_i}\), and
combine the two desired inequalities into the slacks
\[
 F_i(s)=\sigma_iH_i(s)+\delta_iT_k(s),\qquad
 (\sigma_1,\sigma_2)=(-1,1),\qquad \delta_i=\mathbf1_{\{i=1\}}.
\]
Thus we must prove \(F_1,F_2\ge0\) on \([0,\infty)\).  Exact interpolation
already gives
\begin{equation}\label{eq:exact-slack-jets}
 F_i(m)=F_i'(m)=0\qquad(m\in\mathcal N).
\end{equation}
The reference lists are \(X_iq(k)\), extended by zero beyond \(M=168\).
Throughout this section, \(H_i^{\mathrm{ref}}\) denotes the function formed
from those lists with the \emph{exact} direct and transformed spectral
integrals.  Only the finite coefficient rows of Section~\ref{sec:finite} use
quadrature.

The gaps in the full zero set have repeating lengths \(1,2,1,3,2,3\).  The
signed half-gaps centered at the nodes through 88 therefore cover exactly
\[
 [0,s_0],\qquad s_0=89.5.
\]
The last is the right half-gap from 88 to the midpoint 89.5 of the gap to 91.
Every signed half-length satisfies \(1/2\le|y|\le3/2\).  At the origin there
is only the half-gap \([0,1/2]\), with \(\nu=0\); at every positive center,
\(\nu=2\).  Figure~\ref{fig:sign-junction} records the junction with the tail.

\begin{figure}[htbp]
\centering
\begin{tikzpicture}[x=1.15cm,y=1cm,>=Latex,font=\small]
  \draw[->] (-.15,0) -- (8.75,0) node[right] {\(s\)};
  \foreach \a/\b in {1/87,2/88,5/91,7/93} {
    \fill (\a,0) circle (1.8pt);
    \node[below=3pt] at (\a,0) {\(\b\)};
  }
  \draw[densely dashed,gray] (1.5,-.05) -- (1.5,.72);
  \draw[densely dashed,gray] (6,-.88) -- (6,.06);
  \node[above=3pt] at (6,0) {\(92\)};
  \draw[densely dashed] (3.5,-1.05) -- (3.5,1.2);
  \node[above] at (3.5,1.2) {\(s_0=89.5\)};
  \draw[very thick,blue!65!black] (-.1,.65) -- (3.5,.65);
  \draw[blue!65!black] (1.5,.55) -- (1.5,.75);
  \draw[blue!65!black] (2,.55) -- (2,.75);
  \draw[blue!65!black] (3.5,.55) -- (3.5,.75);
  \node[above,text=blue!65!black] at (1.05,.72)
    {finite half-gaps};
  \node[below,text=blue!65!black,font=\scriptsize] at (2.5,.62)
    {center 88};
  \draw[very thick,orange!80!black] (3.5,-.7) -- (8.6,-.7);
  \draw[orange!80!black] (3.5,-.8) -- (3.5,-.6);
  \draw[orange!80!black] (6,-.8) -- (6,-.6);
  \draw[orange!80!black] (8.5,-.8) -- (8.5,-.6);
  \node[below,text=orange!80!black,font=\scriptsize,align=center]
    at (4.75,-.77) {Taylor remainder\\double zero at 91};
  \node[below,text=orange!80!black,font=\scriptsize,align=center]
    at (7.25,-.77) {Taylor remainder\\double zero at 93};
\end{tikzpicture}
\caption{The finite inequalities cover the half-gaps through \(s_0\).
In the tail, 91 is a nearest node on \([89.5,92]\), then 93 on
\([92,94.5]\), and so on.  A double zero at the chosen node turns the uniform
second-derivative estimate into a quadratic remainder bound.}
\label{fig:sign-junction}
\end{figure}

\subsection{Quotients after subtracting values and slopes}

Recall the quotient \(\mathcal D_{m,2}\) introduced before the finite
polynomial rows in Section~\ref{sec:finite}.  For a twice continuously
differentiable function \(u\), Taylor's formula gives
\[
 \mathcal D_{m,2}u(s)
 =\frac{u(s)-u(m)-(s-m)u'(m)}{(s-m)^2}
 =\int_0^1(1-v)u''\bigl(m+v(s-m)\bigr)\,\dd v.
\]
The integral also gives the removable value at \(s=m\).  As before,
\(\mathcal D_{0,0}u=u\) at the origin.  On a prescribed half-gap \(s=m+yx\),
\(0\le x\le1\), the relevant quotient is
\begin{equation}\label{eq:deleted-slack}
 \mathcal D_{m,\nu}F_i(s)
 =\sigma_i\mathcal D_{m,\nu}H_i(s)
  +\delta_i\mathcal D_{m,\nu}T_k(s).
\end{equation}
For \(m>0\), \eqref{eq:exact-slack-jets} makes this equal to
\(F_i(s)/(s-m)^2\) away from the center.  The deletion is performed on each
function before any approximation is made, so a nonzero constant or linear
approximation error is never divided by \((s-m)^2\).

\paragraph{Moving from the exact lists to the reference lists.}
Their combined \(\ell^1\) distance is less than \(2\cdot10^{-8}\).  A unit
list has spectral variation at most 36, direct frequency modulus at most one,
and transformed bounds \(|\lambda|<3\), \(|z|<3\).  The transformed damping
has positive exponent \(\pi\Ima z\), so there is no exponential growth at
nonnegative radii.  Taylor's integral formula and the derivative estimates of
Section~\ref{sec:fourier} therefore give the following bound.  Each
\(H_i-H_i^{\mathrm{ref}}\) contains a direct error from one list and a
transformed error from the other.  The \(\ell^1\) norms of the two list
differences add to less than \(2\cdot10^{-8}\), so the larger of the two
kernel bounds controls the sum:
\begin{align}
 \sup_{s\text{ in the half-gap}}
 \left|\mathcal D_{m,\nu}(H_i-H_i^{\mathrm{ref}})(s)\right|
 &\le\frac{36}{\nu!}(2\cdot10^{-8})
       \max\{\pi^\nu,3(3\pi)^\nu\}\notag\\
 &<.00011.\label{eq:deleted-list-error}
\end{align}
For \(\nu=0\), this is the zeroth-derivative estimate.  For \(\nu=2\), the
factor \(1/2\) is \(\int_0^1(1-v)\,\dd v\), and the last comparison is the
scalar bound in \eqref{eq:scalar-error-budgets}.  No zero jet is required of
the reference function; only the exact slack uses \eqref{eq:exact-slack-jets}.

\paragraph{Replacing the reference quotient by its finite polynomial.}
The entire Taylor series gives, in both cases \(\nu=0,2\),
\[
 \sigma_i\mathcal D_{m,\nu}H_i^{\mathrm{ref}}(m+yx)
 =\sum_{r=0}^{\infty}
   \sigma_i\frac{y^r}{(r+\nu)!}
     (H_i^{\mathrm{ref}})^{(r+\nu)}(m)x^r.
\]
The rows \(l_{i,r}\) in \eqref{eq:bernstein-rows} replace the included
derivatives by quadrature.  We claim
\begin{equation}\label{eq:half-gap-polynomial-error}
 \sup_{0\le x\le1}
 \left|\sigma_i\mathcal D_{m,\nu}H_i^{\mathrm{ref}}(m+yx)
       -\sum_{r=0}^{40}l_{i,r}q(k)x^r\right|<.000001.
\end{equation}
For the included coefficients, Lemma~\ref{lem:quadrature-error} controls the
scaled moment of order \(r+\nu\).  The row above has the additional multiplier
\(y^{-\nu}\), whose absolute value is at most 4.  By
\eqref{eq:reference-norm}, the combined reference-list norm is less than
\(4.8<5\).  Summing the 41 coefficient errors therefore gives at most
\(41\cdot4\cdot5\cdot10^{-31}\).

For the omitted powers, put \(Y=|y|\) and \(l=41+\nu\).  Their total is at
most
\begin{equation}\label{eq:half-gap-series-tail}
 \frac{Y^{-\nu}36\cdot5}
 {l!\bigl(1-3\pi Y/(l+1)\bigr)}
 \left((\pi Y)^l+
 3\max_{1\le j\le6}
 e^{-\pi\Ima z(t_j)m}
       \bigl(\pi|z(t_{j-1})|Y\bigr)^l\right).
\end{equation}
Indeed, each direct derivative contributes a factor \(i\pi t\), and each
transformed derivative a factor \(i\pi z(t)\).  On
\([t_{j-1},t_j]\), both \(\Ima z(t)\) and \(|z(t)|\) decrease.  Thus the
right endpoint maximizes \(e^{-\pi\Ima z(t)m}\), while the left endpoint
bounds the largest frequency modulus, including the endpoint atoms.  For each
fixed frequency, the ratio of successive terms in the absolute
exponential series is at most \(3\pi Y/(l+1)<1\).  Summing this geometric
majorant and then integrating against total variation gives the displayed
denominator.  The factor \(36\cdot5\) bounds the combined spectral
variation of the reference inputs.

The scalar calculation in Appendix~\ref{sec:arithmetic} replaces the maximum
in \eqref{eq:half-gap-series-tail} by the larger sum over all six intervals.
For
\[
 (m,Y,\nu)=(0,1/2,0),\quad(1,1,2),\quad(3,1,2),\quad(4,3/2,2),
\]
it gives, respectively, upper bounds \(1.8\cdot10^{-19}\), \(10^{-11}\),
\(3\cdot10^{-16}\), and \(2\cdot10^{-11}\), all below \(10^{-8}\).  These
four cases cover every half-gap.  The origin is the first case; the largest
half-lengths at 1 and 3 are one; and every later center is at least 4 with
half-length at most \(3/2\).  For fixed \(\nu\), the expression decreases as
\(m\) increases or \(Y\) decreases: its powers contain \(Y^{l-\nu}=Y^{41}\),
and the positive geometric denominator also improves when \(Y\) decreases.
Adding the included quadrature error proves
\eqref{eq:half-gap-polynomial-error}.

\subsection{Lower supports for every target regime}

\begin{lemma}[Exponential supports]\label{lem:exponential-support}
Let \(K,J\) be consecutive endpoints of the parameter partition in
Section~\ref{sec:finite}, and let \(K\le k\le J\), with no upper restriction
when \(J=\infty\).  On every prescribed half-gap, the coefficients \(a_r^*\)
from \eqref{eq:target-support} satisfy
\[
 \mathcal D_{m,\nu}T_k(m+yx)
 \ge\sum_{r=0}^{40}a_r^*x^r-.000001\qquad(0\le x\le1).
\]
\end{lemma}

\begin{proof}
At the origin, \(0\le s\le1/2\), and
\[
 \frac{\partial T_k(s)}{\partial k}
 =T_k(s)\left(1-s-\frac1k\right)>0
\]
because \(k\ge2.36>2\).  Replacing \(k\) by \(K\) therefore lowers the
target.  The Taylor polynomial of \(T_K(yx)\) through degree 40 has
coefficients \((e^K/K)(-Ky)^r/r!\), exactly the first case in
\eqref{eq:target-support}.

At the first positive node, write \(u=s-1\).  The deleted quotient is
\[
 \mathcal D_{1,2}T_k(1+u)=\frac{e^{-ku}-1+ku}{ku^2},
\]
with value \(k/2\) at \(u=0\).  For \(u\ne0\), its derivative with respect to
\(k\) is
\[
 \frac{1-(1+ku)e^{-ku}}{k^2u^2}\ge0,
\]
by \(e^{ku}\ge1+ku\).  The same monotonicity holds at the removable value.
This works for both signed half-gaps, including negative \(u\).  Replacing
\(k\) by \(K\) again lowers the quotient, and its Taylor coefficients are
\(K(-Ky)^r/(r+2)!\), the second case in \eqref{eq:target-support}.

In both cases the left endpoint satisfies \(K\le6\), even for the last,
unbounded parameter interval.  At the origin \(K|y|\le3\), and at node 1
\(K|y|\le6\).  Also \(e^K/K\) increases for \(K>1\), so its value is at most
\(e^6/6\).  The absolute tails after degree 40 are bounded by
\[
 \frac{e^6}{6}\frac{3^{41}}{41!(1-3/42)}<8\cdot10^{-29}<10^{-6},
 \qquad
 \frac{6\cdot6^{41}}{43!(1-6/44)}<10^{-20}<10^{-6},
\]
using the scalar comparisons in Appendix~\ref{sec:arithmetic}.  Thus no
Taylor estimate with unbounded \(k\) is needed.

At every positive node, Taylor's integral formula also gives the general
nonnegative support
\begin{equation}\label{eq:target-integral-quotient}
 \mathcal D_{m,2}T_k(m+u)
 =k\int_0^1(1-v)e^{-k(m-1+vu)}\,\dd v\ge0.
\end{equation}
This handles all cases in which \eqref{eq:target-support} prescribes zero.
It remains to justify the stronger constant at \(m=3\) for finite \(J\).
Here \(-1\le u\le1/2\), so \(a=2+vu\ge1\).  The function \(k e^{-ak}\)
decreases for \(k\ge2.36\), and \eqref{eq:target-integral-quotient} is at least
\[
 J\int_0^1(1-v)e^{-J(2+vu)}\,\dd v.
\]
The probability density \(2(1-v)\) on \([0,1]\) has mean \(1/3\).  Jensen's
inequality for the exponential therefore bounds this expression below by
\[
 \frac J2e^{-J(2+u/3)}\ge\frac J2e^{-13J/6},
\]
which is the third case of \eqref{eq:target-support}.  When \(J=\infty\), the
prescribed support is zero and \eqref{eq:target-integral-quotient} applies.
\end{proof}

\subsection{The bounded radial range}

\begin{lemma}[Signs on the finite half-gaps]\label{lem:half-gap-signs}
For every real \(k\ge2.36\), the exact slacks satisfy \(F_i(s)\ge0\) for
\(0\le s\le s_0\), \(i=1,2\).
\end{lemma}

\begin{proof}
Choose a half-gap and a parameter interval containing \(k\).  Combine the
finite rows and the target support into
\[
 \Phi_i(x;k)=\sum_{r=0}^{40}
       \bigl(l_{i,r}q(k)+\delta_i a_r^*\bigr)x^r.
\]
By the Bernstein conversion defining \eqref{eq:bernstein-rows}, the
coefficient of the \(d\)-th degree-40 Bernstein basis function is
\[
 L_{i,d}q(k)+\delta_i\sum_{r=0}^d\omega_{dr}a_r^*.
\]
The containing-box lower evaluation \eqref{eq:box-lower} applies to the
physical vector \(q(k)\).  Together with
Proposition~\ref{prop:finite-certificate}, it makes every coefficient above
strictly greater than \(.001\).  The Bernstein basis functions are
nonnegative and sum to one on \([0,1]\), so \(\Phi_i(x;k)>.001\) on that
entire interval.

Apply \eqref{eq:deleted-list-error},
\eqref{eq:half-gap-polynomial-error}, and
Lemma~\ref{lem:exponential-support} to \eqref{eq:deleted-slack}.  The result is
the strict quotient margin
\[
 \mathcal D_{m,\nu}F_i(m+yx)
 >.001-.00011-.000001-.000001=.000888>0.
\]
At the origin this is the desired inequality.  At a positive center,
multiplication by \((s-m)^2\) proves the inequality away from the center, and
\eqref{eq:exact-slack-jets} supplies equality at the center.  The half-gaps
cover \([0,s_0]\), while the five parameter intervals cover every finite
\(k\ge2.36\).
\end{proof}

\subsection{The tail and its quadratic scale}

Beyond \(s_0\), the low-node rational part of each direct cardinal function
has a positive signed value.  We will group the high-node terms, the
transformed function, and the target into a single remainder.  That sum has
a small second derivative and exact double zeros; its individual terms need
not have zero jets.  The following geometric bound lets us compare the
positive part and the remainder on the same squared-distance scale.

\begin{lemma}[Quadratic barrier for the sine product]
\label{lem:sine-product-barrier}
For every real \(s\), let \(m\in\mathcal Z\) be any nearest point of the full
periodic zero set, so \(|s-m|=\min_{n\in\mathcal Z}|s-n|\).  Then
\begin{equation}\label{eq:P-quadratic-lower}
 P(s)\ge .68(s-m)^2.
\end{equation}
\end{lemma}

\begin{proof}
The set \(\mathcal Z\) is discrete and periodic, so a nearest point exists.
Between a zero \(m\) and either adjacent midpoint, the quotient
\(P(s)/(s-m)^2\) is positive after its removable value is filled in.  Its
logarithm is concave.  For the factor vanishing at \(m\), with \(u=s-m\),
\[
 \frac{\dd^2}{\dd u^2}
 \log\left(\frac{\sin^2(\pi u/12)}{u^2}\right)
 =\frac2{u^2}-2\left(\frac\pi{12}\right)^2
                \csc^2\left(\frac{\pi u}{12}\right)\le0
\]
for \(u\ne0\), by \(|\sin x|\le|x|\), and the inequality extends through
the removable value.  Every other sine factor has nonpositive logarithmic
second derivative on this half-gap.  A concave function is bounded below by
the smaller of its endpoint values, so the quotient is bounded below by its
smaller value at the center and the midpoint.

At the center, the periodicity of \(P\) reduces the removable value to the
corresponding node value \(Q_n>1.75\) in \eqref{eq:node-constant-bounds}.  At
the midpoints, the exact values are
\[
\begin{array}{c|c}
 \text{gap modulo }12&P(s)/(s-m)^2\text{ at its midpoint}\\ \hline
 {[0,1],\ [3,4]}&12-8\sqrt2\\
 {[1,3]}&1\\
 {[4,7],\ [9,12]}&\dfrac89(3+\sqrt2+\sqrt3+\sqrt6)\\
 {[7,9]}&9
\end{array}
\]
Here the denominator is the square of the half-gap length, independent of
which endpoint is chosen at a tie.  For clarity, the smallest entry can be
derived directly.  At \(s=1/2\), put \(\theta=\pi/24\) and pair the unsquared
factors for residues \((1,7)\), \((3,9)\), and \((0,4)\).  Product-to-sum
gives
\begin{align*}
 (2\sin(-\theta))(2\sin(-13\theta))&=2\sin(\pi/12),\\
 (2\sin(-5\theta))(2\sin(-17\theta))&=2\sin(5\pi/12),\\
 (2\sin\theta)(2\sin(-7\theta))&=1-\sqrt2.
\end{align*}
The product of the first two right sides is one.  Thus
\(P(1/2)/(1/2)^2=4(1-\sqrt2)^2=12-8\sqrt2\).  The same product-to-sum
calculation at the other midpoints gives the remaining entries.  The identity
\(P(4-s)=P(s)\), which follows from invariance of \(L\) under
\(a\mapsto4-a\) modulo 12, accounts for the paired gap types in the table.
The smallest entry satisfies \(12-8\sqrt2>.6862>.68\), as recorded in
Appendix~\ref{sec:arithmetic}; every other entry is larger than \(.68\) as
well.  The half-gaps cover the real line, proving the claim, including the
zero values themselves.
\end{proof}

\begin{lemma}[Signs in the tail]\label{lem:tail-signs}
For every real \(k\ge2.36\), the exact slacks satisfy \(F_i(s)\ge0\) for
\(s\ge s_0\), \(i=1,2\).
\end{lemma}

\begin{proof}
We first bound the low-node rational contribution, then compare the remaining
term on the same quadratic scale.

\paragraph{The low-node rational part.}
Separate each direct cardinal function as
\begin{align*}
 R_{i,0}(s)
 &=C_i+\sum_{\substack{n\in\mathcal N_0\\n\le88}}
       \left(\frac{c_{in}}{(s-n)^2}+\frac{d_{in}}{s-n}\right),\\
 p_i(s)&=P(s)R_{i,0}(s)+p_{i,\mathrm{high}}(s).
\end{align*}
The high function contains precisely the node columns beginning at 91 and no
constant column.  Every denominator in \(R_{i,0}\) is positive for
\(s\ge s_0\).  We claim
\begin{equation}\label{eq:tail-rational-positive}
 \sigma_iR_{i,0}(s)>.002\qquad(s\ge s_0).
\end{equation}
For the reference lists, use the decomposition preceding
\eqref{eq:tail-certificate}, which groups the simple-pole rows into their
common \(1/s\) numerator before taking negative parts.  In the parameter
box containing \(k\), each numerator row \(l\) has the nonpositive lower
bound \(\eta_{K,J}(l)=\min(0,[l]_{K,J})\).  The reciprocal factors are
nonnegative and nonincreasing for \(s\ge s_0\), so their products with these
lower bounds give a value
at least as large as at \(s_0\); the constant row retains its full lower
bound \([C_i^*]_{K,J}\).  Thus \eqref{eq:tail-certificate} gives
\(\sigma_iR_{i,0}^{\mathrm{ref}}(s)>.0021\).

The passage to the exact list changes this rational function by less than
\(2\cdot10^{-8}\).  Indeed, \(s-n\ge1.5\) for every low node, so both
reciprocal coefficient factors have modulus at most one, as does the constant
factor; apply the full-list correction in \eqref{eq:exact-list-bounds}.
Since \(.0021-2\cdot10^{-8}>.002\), this proves
\eqref{eq:tail-rational-positive}.

\paragraph{The high-node list and the remainder.}
For the exact pair define the coupled input sizes
\[
 w=\sum_{i=1}^2\sum_{n\in\mathcal N_0}(|c_{in}|+|d_{in}|),\qquad
 c=\sum_{i=1}^2|C_i|.
\]
The full-list norm and the fixed extras give \(w+c<5\) and \(c<.02\), by
\eqref{eq:exact-list-bounds} and \eqref{eq:extra-constant-bound}.  At every
node \(n\ge91\), the exact equations
\eqref{eq:infinite-interpolation-system} read
\[
 (c_{1n},d_{1n})=(S\ell_2)_n+(\mathcal JT_k)_n,
 \qquad
 (c_{2n},d_{2n})=(S\ell_1)_n.
\]
Here \(S\ell_i=-\mathcal JK_{p_i}\) includes all input coordinates,
including the extras.  Applying Lemma~\ref{lem:kernel-bounds} to both lists
with their summed input bounds \(w,c\), the combined high-node norm is at
most \(U_{91}(w,c)\) plus the target geometric majorant
\[
 \frac{1+(1+2.12)/2.36}{1.75}
       \frac{e^{-2.36\cdot90}}{1-e^{-2.36}}<10^{-89}.
\]
This follows from \eqref{eq:target-jet-bound} by overcounting with all integers
from 91; the last comparison is the local calculation
\eqref{eq:target-tail-91}.  The finer scalar bound in
\eqref{eq:scalar-error-budgets} is \(U_{91}(5,.02)<2.2711\cdot10^{-7}\).
Since \(U_{91}\) is increasing in both input bounds, we obtain
\[
 \sum_{i=1}^2\sum_{\substack{n\in\mathcal N\\n\ge91}}
       (|c_{in}|+|d_{in}|)
 <2.2711\cdot10^{-7}+10^{-89}<2.4\cdot10^{-7}.
\]
The difference between \(2.2711\cdot10^{-7}\) and
\(2.4\cdot10^{-7}\) leaves room for this target contribution.

Define
\[
 Z_i(s)=\sigma_i\bigl(p_{i,\mathrm{high}}(s)+K_{p_{3-i}}(s)\bigr)
                  +\delta_iT_k(s).
\]
Then
\begin{equation}\label{eq:tail-decomposition}
 F_i(s)=\sigma_iP(s)R_{i,0}(s)+Z_i(s).
\end{equation}
We have the uniform bound
\begin{equation}\label{eq:tail-second-derivative}
 \sup_{s\ge s_0}|Z_i''(s)|<.0002.
\end{equation}
To verify it, the direct high columns contribute at most
\(36\pi^2(2.4\cdot10^{-7})\).  For the low-frequency transformed term, the
actual coupled full-list norm gives the important estimate
\[
 25c+36w\le36(c+w)<180.
\]
Thus the \([0,5/6]\) variation in \eqref{eq:refined-derivative-bound}, which
includes the atom at \(5/6\), is bounded by \(36\cdot5\).  On \((5/6,1]\),
we may separately use \(w<5\) and \(c<.02\); the only atom there is at 1 and
contributes \(2c<.04\).  Finally, \(T_k''(s)=k e^{-k(s-1)}\).  Its supremum
over \(s\ge s_0\), \(k\ge2.36\), occurs at \(s=s_0\), \(k=2.36\), since
\(s_0-1>1/2.36\).  The resulting four-term bound is
\begin{align*}
 &36\pi^2(2.4\cdot10^{-7})
  +3(3\pi)^2\,36\cdot5\,e^{-.70s_0}\\
 &\quad+1.25(1.32\pi)^2
       \left(\frac{2.54\cdot5}{s_0}+.04\right)e^{-.187s_0}
  +2.36e^{-2.36(s_0-1)}
 <8.5485\cdot10^{-5}<.0002,
\end{align*}
by the scalar comparison in Appendix~\ref{sec:arithmetic}.  This proves
\eqref{eq:tail-second-derivative}.

At every node \(m\in\mathcal N\) with \(m\ge91\), \(R_{i,0}\) is regular and
\(P\) has a double zero.  Subtracting its product from the exact slack jets
in \eqref{eq:tail-decomposition} therefore gives
\begin{equation}\label{eq:tail-remainder-jets}
 Z_i(m)=Z_i'(m)=0\qquad(m\in\mathcal N,\ m\ge91).
\end{equation}
These zeros use the exact interpolation equations, not merely the small norm
of the high list.

\paragraph{Comparison at a nearest zero.}
For \(s\ge s_0\), choose a nearest point \(m\in\mathcal Z\) with \(m\ge91\).
Such a choice is possible throughout the tail; at \(s=s_0\), choose 91 from
the tied points 88 and 91.  The segment joining \(s\) and \(m\) stays in
\([s_0,\infty)\).  Taylor's integral formula, using
\eqref{eq:tail-second-derivative} and \eqref{eq:tail-remainder-jets}, gives
\[
 |Z_i(s)|\le .0001(s-m)^2.
\]
On the other hand, \eqref{eq:tail-rational-positive} and
Lemma~\ref{lem:sine-product-barrier} imply
\[
 \sigma_iP(s)R_{i,0}(s)\ge .68\cdot.002(s-m)^2.
\]
Since \(.68\cdot.002>.0001\), \eqref{eq:tail-decomposition} gives
\(F_i(s)\ge0\).  At a node the exact zero value gives the same conclusion,
without division by a vanishing factor.
\end{proof}

\begin{proof}[Proof of Theorem~\ref{thm:normalized-pair}]
For every real \(k\ge2.36\), Proposition~\ref{prop:exact-interpolation}
supplies real summable lists with the prescribed extras.  Section~\ref{sec:fourier}
makes \(H_1,H_2\) entire and real on the real axis, and
Lemma~\ref{lem:fourier-pair} gives the real radial Schwartz Fourier pair
\[
 f(x)=e^{-\pi hb|x|^2}H_1(b|x|^2),\qquad
 \widehat f(\xi)=e^{-\pi hb|\xi|^2}H_2(b|\xi|^2).
\]
Lemmas~\ref{lem:half-gap-signs} and \ref{lem:tail-signs} cover all \(s\ge0\),
so \(H_1(s)\le T_k(s)\) and \(H_2(s)\ge0\).  The exact jets are
\[
 H_1(n)=T_k(n),\qquad H_1'(n)=T_k'(n),\qquad
 H_2(n)=H_2'(n)=0\qquad(n\in\mathcal N).
\]
This proves the theorem.
\end{proof}
